%%%%%%%%%%%%%%%%%%%%%%%%%%%%%%%%%%%%%%%%%%%%%%%%%%%%%
%%                 1.1 JUSTIFICATIVA               %%
%%%%%%%%%%%%%%%%%%%%%%%%%%%%%%%%%%%%%%%%%%%%%%%%%%%%%

\section{Justificativa}
\label{sec:justificativa}

A realização domiciliar de exercícios fisioterapêuticos depende, em grande parte, da capacidade do paciente de manter a regularidade e de executar as atividades conforme as orientações recebidas. Fatores pessoais, características do tratamento e dificuldades práticas podem comprometer a adesão a essas atividades \cite{Jack2010}. Para pessoas idosas, essas barreiras podem se tornar mais relevantes quando existem limitações de visão, audição, mobilidade, atenção ou pouca familiaridade com interfaces digitais. Nesse cenário, justifica-se investigar ferramentas que ofereçam apoio imediato entre as sessões presenciais, sem substituir o acompanhamento profissional.

Do ponto de vista técnico, a estimativa de pose permite obter pontos corporais a partir de câmeras convencionais, dispensando sensores vestíveis, marcadores físicos ou câmeras de profundidade \cite{debnath2022}. Soluções como o BlazePose foram projetadas para execução em tempo real em dispositivos móveis \cite{Bazarevsky2020}, e o uso do MediaPipe já foi avaliado em exercícios fisioterapêuticos \cite{simoes2024}. Isso torna plausível a construção de uma solução de baixo custo, baseada em equipamentos que o usuário já possui, como um smartphone.

Entretanto, detectar o movimento não é suficiente. A literatura aponta desafios na geração de feedback a partir da estimativa de pose \cite{tharatipyakul2024}, e a qualidade da detecção pode ser afetada por iluminação, enquadramento, oclusão de articulações e distância da câmera. Assim, um sistema que classifique movimentos apenas por valores angulares fixos corre o risco de apresentar um problema de captura como se fosse um erro do paciente, ou de desconsiderar as diferenças individuais de amplitude. A contribuição técnica pretendida pelo FisioMove está justamente na integração entre três aspectos: a calibração individual, a verificação da confiabilidade dos pontos corporais antes da emissão de feedback e o uso exclusivo de parâmetros definidos pelo fisioterapeuta responsável.

A escolha do recorte também se justifica pela viabilidade de um projeto acadêmico. O protótipo será limitado a dois exercícios, o que permite analisar cada movimento com maior profundidade e submeter suas regras à validação de um profissional de Fisioterapia, em vez de oferecer um catálogo amplo sem avaliação adequada. Pelo mesmo motivo, o sistema não realiza diagnóstico, não prescreve tratamentos e não altera automaticamente metas ou amplitudes prescritas.

A gamificação foi incorporada como recurso de apoio ao engajamento, considerando que revisões sobre o tema em aplicações de saúde indicam seu potencial, desde que os mecanismos sejam adequados à finalidade e ao público da aplicação \cite{Sardi2017}. No contexto da reabilitação musculoesquelética, esses recursos também têm sido investigados como estratégia complementar \cite{alfieri2022}. No FisioMove, os elementos lúdicos serão vinculados ao cumprimento do plano prescrito, e não à realização de repetições além do recomendado, evitando incentivar esforço excessivo. Como o público prioritário é composto por pessoas idosas, a interface deve ser minimalista, com poucas informações e poucos controles por tela, o que orienta a seleção de mecanismos simples de compreender.

Por fim, a proposta considera a privacidade dos usuários desde a concepção. Como imagens de câmera e informações relacionadas à saúde são dados sensíveis \cite{LGPD2018}, a arquitetura adotada processa a estimativa de pose no próprio dispositivo e envia ao servidor somente dados derivados das sessões, como quantidade de repetições, duração e conclusão das metas. Essa decisão reduz a exposição de imagens, a dependência de conexão e o custo de infraestrutura, contribuindo para uma solução mais segura e adequada ao contexto domiciliar.

Dessa forma, o trabalho se justifica por sua relevância social, ao buscar apoiar a continuidade de exercícios domiciliares por pessoas idosas; por sua relevância técnica, ao investigar o uso responsável da estimativa de pose em dispositivos convencionais; e por sua relevância acadêmica, ao discutir as possibilidades e os limites dessa abordagem sem atribuir ao sistema afirmações clínicas que não possam ser sustentadas.
